\documentclass[../main.tex]{subfiles}
\begin{document}

\section{Latent analysis}
\label{sec:latent}

\begin{claim}{STRONG}{F12}
If the backdoor in a ViT is a single linear direction in the residual stream, written late and routed through attention from the trigger's patches to the class token, then removing that direction alone should kill the attack while zeroing any set of coordinates should not, the direction's alignment to its final form should assemble over a band of layers rather than persist from the start, and patching the trigger's own tokens with their clean values should restore the clean answer fully before the routing depth and not after. Each is tested below, and each states whether it rests on a single checkpoint or on a bootstrap over the panel.
\end{claim}

\subsection{One direction, not a set of neurons}

\Cref{tab:direction-ablation} is the causal test, on \AblationCheckpoints{} CIFAR-10 checkpoints at the highest panel rate, a checkpoint per attack\finding{STRONG}{F12}. Projecting out the lowest-rank backdoor direction after the final LayerNorm, where the head reads, leaves attack success at most \AblationDirectionMaxAsr{} on the all-to-one attacks, while a random direction leaves it at \AblationRandomDirectionMinAsr{} at the lowest and zeroing the most trigger-activated coordinates leaves it at \AblationTopCoordinatesMinAsr{} at the lowest. The direction is real, causal and rotated off the coordinate basis, which is why every earlier attempt of this project to exploit units, heads and coordinate rankings failed and why a structured mask cannot beat unstructured noise on the operator axis. \Cref{fig:direction} gives the depth profile on the same checkpoints, relative direction norm, mean trigger-activated change and clean-versus-triggered CKA per layer, with the benign reference probed by the same trigger peaking at \DirectionBenignPeakRelNorm{} against a smallest final value of \DirectionBackdooredMinFinalRelNorm{} on the backdoored ones. This is a single dataset and a single checkpoint per attack, with a benign control and the ablation's own controls.

% generated by scripts/paper/mech_direction.py from results/backdoor_neuron_ablation.json at 2026-09-23T11:19:21.201473+00:00, commit fe1619cdcc6be48ee0458803c6a7dd41e209c3bc-dirty. Do not edit by hand.
\begin{table}[htbp]
\centering
\small
\caption{The causal ablation at the peak layer, CIFAR-10 at the highest panel rate: attack success over clean accuracy after projecting out the rank-1 backdoor direction, a random direction, the top-20 trigger-activated coordinates and 20 random coordinates, as a forward hook after the final LayerNorm.}
\label{tab:direction-ablation}
\begin{adjustbox}{max width=\linewidth}
\begin{tabular}{lrrrrrrr}
\toprule
attack & peak layer & ASR & CA & rank-1 direction & random direction & top-20 coordinates & random-20 coordinates \\
\midrule
\bottomrule
\end{tabular}
\end{adjustbox}
\end{table}


\begin{figure}[t]
\centering
\includegraphics[width=\linewidth]{figures/mech_direction.pdf}
\caption{Depth profile of the backdoor direction on the CIFAR-10 checkpoints, benign reference dashed. From \code{paper/figures/mech\_direction.json}.}
\label{fig:direction}
\end{figure}

\subsection{Crystallisation and the per-layer TAC pass}

The direction does not persist uniformly\finding{SUPPORTING}{F13}. \Cref{tab:direction-crystallisation} reads, on \CrystallisationCheckpoints{} CIFAR-100 and Tiny checkpoints across every poison rate trained, the first layer at which a layer's direction aligns at least half with the final one. It is \CrystallisationBadnetLayer{} for BadNets and \CrystallisationBlendLayer{} for Blend, with the poison rate moving neither, so the global trigger crystallises earlier than the patch. \Cref{tab:tac-layers} and \cref{fig:tac-layers} extend the depth profile to \TacLayersCheckpoints{} checkpoints on GTSRB, CIFAR-100 and Tiny with a GPU pass run for this paper on \TacLayersSamples{} paired images each, and the onset layer, where the relative norm first reaches half its final value, is \TacLayersBadnetOnset{} for BadNets and \TacLayersBlendOnset{} for Blend on average, with the largest benign peak at \TacLayersBenignPeakMax{}. The earlier CIFAR-10 ordering, distributed triggers onset early and the patch trigger late, holds on CIFAR-100 and Tiny and reverses on GTSRB, where Blend's onset is later than BadNets', so the onset ordering is a property of the dataset as well as of the trigger and the per-attack band rule that was built on it is not recoverable from this pass either.

% generated by scripts/paper/mech_direction.py from results/direction_persistence.json at 2026-09-11T04:08:14.450470+00:00, commit f2d5782d3d0417e887985131455a9b38b72c2b27-dirty. Do not edit by hand.
\begin{table}[htbp]
\centering
\small
\caption{Crystallisation depth on CIFAR-100 and Tiny across poison rates: the first layer whose backdoor direction aligns at cosine at least 0.5 with the final layer's, per checkpoint, summarised per attack.}
\label{tab:direction-crystallisation}
\begin{adjustbox}{max width=\linewidth}
\begin{tabular}{lrrrr}
\toprule
attack & checkpoints & mean layer & earliest & latest \\
\midrule
BadNets & 8 & 10.4 & 10 & 11 \\
Blend & 8 & 8.2 & 8 & 9 \\
\bottomrule
\end{tabular}
\end{adjustbox}
\end{table}

% generated by scripts/paper/mech_tac_layers.py from results/_experiments/tac_layers/tac_layers.json at 2026-09-23T18:21:02.095010+00:00, commit 5a1f1c5814dc5676c093f3021490a849712a13f4-dirty. Do not edit by hand.
\begin{table}[htbp]
\centering
\small
\caption{Per-layer profile of the backdoor direction on 18 ViT-B/16 models, 500 paired images each.}
\label{tab:tac-layers}
\begin{adjustbox}{max width=\linewidth}
\begin{tabular}{llrrrrr}
\toprule
Dataset & Model & Peak layer & Peak rel. norm & Onset layer & Final CKA & Final TAC \\
\midrule
GTSRB & BadNets 10\% & 12 & 1.64 & 7 & 0.57 & 1.05 \\
GTSRB & benign & 5 & 0.01 & 2 & 1.00 & 0.01 \\
GTSRB & Blend 10\% & 12 & 3.40 & 9 & 0.12 & 0.62 \\
GTSRB & BPP 10\% & 9 & 1.66 & 7 & 0.41 & 0.76 \\
GTSRB & LF 10\% & 12 & 1.68 & 7 & 0.48 & 0.86 \\
GTSRB & TaCT 10\% & 12 & 0.38 & 12 & 0.99 & 0.20 \\
GTSRB & WaNet 10\% & 12 & 1.38 & 8 & 0.29 & 0.59 \\
CIFAR-100 & BadNets 1\% & 12 & 0.79 & 8 & 0.44 & 0.57 \\
CIFAR-100 & BadNets 10\% & 10 & 0.92 & 8 & 0.39 & 0.57 \\
CIFAR-100 & benign & 10 & 0.12 & 6 & 0.99 & 0.06 \\
CIFAR-100 & Blend 1\% & 10 & 1.04 & 7 & 0.23 & 0.62 \\
CIFAR-100 & Blend 10\% & 9 & 1.70 & 5 & 0.16 & 0.71 \\
CIFAR-100 & BPP 10\% & 12 & 1.28 & 8 & 0.20 & 0.75 \\
CIFAR-100 & LF 10\% & 12 & 1.46 & 8 & 0.24 & 0.75 \\
CIFAR-100 & TaCT 10\% & 10 & 0.26 & 9 & 0.98 & 0.13 \\
Tiny ImageNet & BadNets 10\% & 12 & 0.99 & 9 & 0.41 & 0.55 \\
Tiny ImageNet & benign & 9 & 0.04 & 5 & 1.00 & 0.02 \\
Tiny ImageNet & Blend 10\% & 10 & 1.52 & 7 & 0.15 & 0.66 \\
\bottomrule
\end{tabular}
\end{adjustbox}
\end{table}


\begin{figure}[t]
\centering
\includegraphics[width=\linewidth]{figures/mech_tac_layers.pdf}
\caption{Relative backdoor-direction norm and mean trigger-activated change by layer on GTSRB, CIFAR-100 and Tiny, benign references dashed. From \code{paper/figures/mech\_tac\_layers.json}.}
\label{fig:tac-layers}
\end{figure}

\subsection{Activation patching}

\Cref{tab:activation-patching} and \cref{fig:patching} average, over every cached cell of each attack, the recovery of the clean answer when a token group at a layer of the residual stream is overwritten with its value on the same image without the trigger, the paired counterfactual \citet{zhang2024patching} recommend over noise corruption\finding{STRONG}{F14}. For the local triggers, BadNets and TaCT, patching the trigger's own tokens restores the clean answer at \PatchingLocalEarlyTrigger{} over layers \PatchingEarlyLayers{}, the tokens alone still suffice at layer \PatchingBadnetLastFull{} for BadNets, the class token first recovers more than the trigger tokens at layer \PatchingBadnetCrossover{}, and a random token group of the same count recovers at most \PatchingRandomNullMax{} anywhere on the global attacks and at most a few thousandths on the local ones. The null is the qualification. A global trigger covers every patch, so its trigger set is all \VitPatches{} tokens and a same-count random group is the same set, which is why the global rows read the same on the null and why the patching result is a local-trigger result on \PatchingCells{} cells over \PanelDatasetsCached{} datasets. On those, the backdoor is causally located entirely in its own patches through the first half of the network and reaches the class token only in the last quarter.

% generated by scripts/paper/mech_activation_patching.py from results/coverage/coverage.json, results/<folder>/activation_patching.json (65 cells) at 2026-09-23T14:31:39.049124+00:00, commit 039fc6fea9212f8f6d49218ed14b336aa8f6a0d3-dirty. Do not edit by hand.
\begin{table}[htbp]
\centering
\small
\caption{Activation patching at the residual stream, averaged over the models whose attack succeeded, per attack: mean recovery when the trigger's own tokens are patched over the first 6 layers, the last layer at which those tokens alone recover at least half the clean answer, the first layer at which the class token recovers more than the trigger tokens, the class token's recovery at the last layer, and the largest recovery a random token group of the same count reaches at any layer, the null.}
\label{tab:activation-patching}
\begin{adjustbox}{max width=\linewidth}
\begin{tabular}{llrrrrrr}
\toprule
attack & trigger & n & trigger tokens, layers 1 to 6 & last layer trigger tokens suffice & crossover layer & class token, last layer & random null, max \\
\midrule
BPP & global & 12 & 0.993 & 10 & 11 & 0.656 & 1.000 \\
WaNet & global & 5 & 0.992 & 7 & 8 & 0.918 & 1.000 \\
TaCT & local & 11 & 1.000 & 9 & 11 & 0.545 & 0.001 \\
SIG & global & 1 & 0.993 & 11 & 12 & 0.744 & 1.000 \\
BadNets & local & 12 & 1.000 & 10 & 12 & 0.661 & 0.005 \\
Blend & global & 12 & 0.934 & 8 & 11 & 0.755 & 0.966 \\
LF & global & 12 & 0.996 & 9 & 10 & 0.782 & 1.000 \\
\bottomrule
\end{tabular}
\end{adjustbox}
\end{table}


\begin{figure}[t]
\centering
\includegraphics[width=\linewidth]{figures/mech_activation_patching.pdf}
\caption{Activation patching at the residual stream, mean recovery over cached cells per attack. From \code{paper/figures/mech\_activation\_patching.json}.}
\label{fig:patching}
\end{figure}

\subsection{Routing, sinks and artifact tokens}

\Cref{tab:routing} and \cref{fig:routing} consolidate the routing records\finding{SUPPORTING}{F15}. On triggered images the class token's attention mass on BadNets' trigger tokens is \RoutingBadnetClsEarly{} over blocks \BandEarly{}, \RoutingBadnetClsRouting{} over blocks \BandMiddle{} and \RoutingBadnetClsLate{} over blocks \BandLate{}, against \RoutingBenignClsLate{} late on the benign reference, and the global triggers read \RoutingGlobalClsRouting{} on the middle band because their trigger set is every patch. Attention's share of the summed target-logit attribution of the sublayers' writes is \RoutingBadnetAttentionShare{} for BadNets, so on the logit the attention write dominates for the patch trigger, and the MLP being token-wise cannot move content between tokens whatever it writes. The trigger's tokens carry a norm rank of \RoutingLocalRankRouting{} among patches over the middle band for the local triggers, against \RoutingBenignRankRouting{} on the benign reference under the same trigger, which is the raw pixel patch still being visible there and is why the first blocks are excluded from every norm statistic, and the peak max-over-median patch-norm ratio from the third block on is \RoutingBadnetRatioPeak{} for BadNets against \RoutingBenignRatioPeak{} on clean images. When the trigger's tokens are masked the class token's attention mass is restored by other tokens at a takeover fraction of \RoutingBadnetTakeover{} over the middle band for BadNets against \RoutingBenignTakeover{} on the benign reference, so the trigger's role is one other tokens step into only partly. This ViT has no native register population, the local trigger manufactures one and the global triggers do not, and the artifact-token detector built on the norm ratio is a single-cell observation under audit rather than a result.

% generated by scripts/paper/mech_routing.py from results/coverage/coverage.json, results/<folder>/cls_routing.json (30), results/<folder>/logit_attribution.json (30), results/<folder>/artifact_tokens.json (69), results/<folder>/sink_anatomy.json (69) at 2026-09-23T14:35:23.731709+00:00, commit 039fc6fea9212f8f6d49218ed14b336aa8f6a0d3-dirty. Do not edit by hand.
\begin{table}[htbp]
\centering
\small
\caption{Routing and sink anatomy per attack. CLS attention mass on the trigger tokens of triggered images, averaged over the early, routing and late layer bands (cls\_routing.json, highest panel rate). Attention's share of the summed target-logit attribution of the 2 sublayers' writes (logit\_attribution.json). The trigger tokens' norm rank among patches over the routing and late bands, the peak max-over-median patch-norm ratio from layer 3 on (artifact\_tokens.json, every model whose attack succeeded), and the takeover fraction over the routing band, how much attention mass other tokens absorb when the trigger's are masked (sink\_anatomy.json). A global trigger covers every patch, so its token-level columns are absent. n is routing/artifact models.}
\label{tab:routing}
\begin{adjustbox}{max width=\linewidth}
\begin{tabular}{llrrrrrrrrr}
\toprule
attack & trigger & n & CLS mass 1-4 & CLS mass 5-8 & CLS mass 9-12 & attn share & rank 5-8 & rank 9-12 & peak ratio & takeover 5-8 \\
\midrule
BPP & global & 4/12 & 0.850 & 0.825 & 0.968 & 0.87 & -- & -- & 3.51 & -- \\
WaNet & global & 3/5 & 0.840 & 0.824 & 0.938 & 0.55 & -- & -- & 3.65 & -- \\
TaCT & local & 4/11 & 0.019 & 0.037 & 0.157 & 0.56 & 0.92 & 0.61 & 4.06 & 0.94 \\
SIG & global & 1/1 & 0.956 & 0.911 & 0.984 & 0.67 & -- & -- & 1.65 & -- \\
Label-Consistent & local & 2/0 & 0.102 & 0.191 & 0.511 & 0.84 & -- & -- & -- & -- \\
BadNets & local & 4/12 & 0.025 & 0.231 & 0.550 & 0.83 & 0.89 & 0.58 & 4.26 & 0.50 \\
Blend & global & 4/12 & 0.850 & 0.896 & 0.967 & 0.54 & 0.00 & 0.00 & 3.32 & 4.28 \\
LF & global & 4/12 & 0.867 & 0.856 & 0.971 & 0.55 & -- & -- & 2.74 & -- \\
benign & control & 4/4 & 0.020 & 0.029 & 0.025 & 0.78 & 0.91 & 0.68 & 3.65 & 0.95 \\
\bottomrule
\end{tabular}
\end{adjustbox}
\end{table}


\begin{figure}[t]
\centering
\includegraphics[width=\linewidth]{figures/mech_routing.pdf}
\caption{Left, CLS attention mass on the trigger tokens by layer per attack at the highest panel rate. Right, the trigger tokens' norm rank among patches by layer over the cached cells. Benign references dashed. From \code{paper/figures/mech\_routing.json}.}
\label{fig:routing}
\end{figure}

\subsection{Prediction depth}

Prediction depth \citep{baldock2021difficulty} reads every block through the network's own final LayerNorm and head and averages the lens probability of the final answer over depth, a single forward pass and no rate to pick\finding{STRONG}{F17}. \Cref{tab:prediction-depth} pairs it against the recommended placement on all \DepthCells{} cached cells. Its AUROC is \DepthAurocAll{}, a paired delta of \DepthMinusPsbdAuroc{} with interval \DepthMinusPsbdAurocCi{}, and its TPR at the \LowestQuantilePercent{} budget is \DepthMinusPsbdTpr{} above PSBD's with interval \DepthMinusPsbdTprCi{}, winning on \DepthTprWins{} cells. The gain is concentrated on the hard attacks, \DepthMinusPsbdTprHard{} at the low budget with interval \DepthMinusPsbdTprHardCi{}, and it loses on BadNets at every rate, so the earlier reading of a large low-budget advantage held on the subset it was read on and does not on the panel. The benign references read \DepthBenignAuroc{}. A statistic with no perturbation at all reaching about the same AUROC as the best perturbation says that what PSBD measures is a property of the sample's trajectory through depth as much as of its robustness, and the pair are complementary at the operating point PSBD is weakest at.

% generated by scripts/paper/mech_prediction_depth.py from /lustre/home/pstika/projects/PSBD-ViT/results/coverage/coverage.json, /lustre/home/pstika/projects/PSBD-ViT/results/<folder>/prediction_depth.json (65 cells), /lustre/home/pstika/projects/PSBD-ViT/results/<folder>/psbd_metrics.json at 2026-09-11T02:52:25.929444+00:00, commit 8e51d579049371c06dcaf4c318714f3b070c9656-dirty. Do not edit by hand.
\begin{table}[t]
\centering
\small
\caption{Prediction depth (depth\_soft, 1 forward pass) against PSBD at the recommended placement and the adaptive rule, paired within cell on the same splits: AUROC at the headline quantile and TPR at the smallest budget, with the paired delta, its bootstrap interval and the count of cells where prediction depth wins at the low budget.}
\label{tab:prediction-depth}
\begin{tabular}{lrrrlrrlr}
\toprule
subset & n & depth AUROC & PSBD AUROC & delta AUROC [CI] & depth TPR@1\% & PSBD TPR@1\% & delta TPR [CI] & depth wins \\
\midrule
all & 65 & 0.911 & 0.935 & -0.024 [-0.078, +0.023] & 0.648 & 0.629 & +0.019 [-0.079, +0.113] & 41 \\
hard attacks & 29 & 0.883 & 0.876 & +0.007 [-0.102, +0.094] & 0.553 & 0.371 & +0.183 [+0.062, +0.296] & 19 \\
easy attacks & 36 & 0.934 & 0.983 & -0.049 [-0.090, -0.015] & 0.724 & 0.836 & -0.112 [-0.237, +0.009] & 22 \\
bpp & 12 & 0.977 & 0.948 & +0.029 [-0.001, +0.068] & 0.898 & 0.757 & +0.141 [-0.046, +0.318] & 10 \\
wanet & 5 & 0.948 & 0.845 & +0.104 [+0.009, +0.276] & 0.895 & 0.290 & +0.605 [+0.353, +0.807] & 5 \\
tact & 11 & 0.754 & 0.853 & -0.099 [-0.361, +0.102] & 0.059 & 0.018 & +0.041 [-0.028, +0.155] & 3 \\
sig & 1 & 0.843 & 0.418 & +0.425 -- & 0.157 & 0.027 & +0.129 -- & 1 \\
badnet\_a2o & 12 & 0.845 & 0.992 & -0.146 [-0.236, -0.064] & 0.421 & 0.776 & -0.355 [-0.626, -0.076] & 5 \\
blend & 12 & 0.976 & 0.978 & -0.002 [-0.030, +0.020] & 0.833 & 0.829 & +0.003 [-0.162, +0.153] & 6 \\
lf & 12 & 0.981 & 0.980 & +0.001 [-0.003, +0.005] & 0.919 & 0.904 & +0.015 [-0.083, +0.087] & 11 \\
rate 0.01 & 19 & 0.846 & 0.944 & -0.098 [-0.227, -0.002] & 0.486 & 0.635 & -0.148 [-0.347, +0.046] & 8 \\
rate 0.05 & 22 & 0.916 & 0.954 & -0.038 [-0.132, +0.030] & 0.666 & 0.661 & +0.005 [-0.154, +0.152] & 14 \\
rate 0.1 & 24 & 0.959 & 0.911 & +0.048 [+0.001, +0.105] & 0.760 & 0.594 & +0.166 [+0.050, +0.289] & 19 \\
\bottomrule
\end{tabular}
\end{table}


\end{document}
